\documentclass{article}

\usepackage{makecell, graphicx}
\graphicspath{ {./images/} }

\title{Ch. 15 CW \#1 Solutions}
\author{Principles of Microeconomics}
\date{\today}

\begin{document}

\maketitle

\section*{15.1. Warm Up}

\begin{enumerate}

\item The market needs to have many buyers and many sellers, and the goods offered by every seller need to be largely the same.

\item We assume that firms are trying to maximize their profit. 

\end{enumerate}

\section*{15.1. What is a Competitive Market}

\begin{enumerate}

	
\item (c) -- In a perfectly competitive market, firms are price takers.

\item

	\begin{enumerate}
	
	\item Milk is essentially the same regardless of the producer, and there are many buyers and sellers.
	
	\item
	
		\begin{center}
		\begin{tabular}{c | c | c | c | c }
		Quantity & Price & \makecell{Total\\Revenue} & \makecell{Average\\Revenue} & \makecell{Marginal\\Revenue} \\
		\hline
		1 gallon & \$6 & \$6 & \$6 & \$6 \\
		\hline
		2 & 6 & 12 & 6 & 6 \\
		\hline
		3 & 6 & 18 & 6 & 6 \\
		\hline
		4 & 6 & 24 & 6 & 6 \\
		\hline
		5 & 6 & 30 & 6 & 6 \\
		\hline
		6 & 6 & 36 & 6 & 6 \\
		\hline
		7 & 6 & 42 & 6 & 6 \\
		\hline
		8 & 6 & 48 & 6 & 6
		\end{tabular}
		\end{center}
	
	\end{enumerate}
	
\item (c) -- Marginal revenue is the price of the good which is constant in a perfectly competitive industry. Total revenue is $PQ$, so it's proportional to the quantity sold. If quantity increases by 10\%, total revenue increases by 10\%.

\end{enumerate}

\end{document}